\chapter*{Kurzfassung} 
\markboth{Kurzfassung}{Kurzfassung}
%% ==============================
Die Kurzfassung (Englisch abstract) hat das Ziel, den Leser über die wesentlichen Aspekte der Arbeit zu informieren. Hierzu gehören die Problemstellung, die Zielsetzung, der Lösungsweg sowie die wichtigsten Erkenntnisse der Arbeit. Dies alles soll in kurzer prägnanter Form auf maximal einer DIN A4-Seite zusammengefasst werden und keine Formeln oder Literaturhinweise enthalten. Dem Leser soll schnell erkennbar sein, was ihn beim Lesen der Arbeit erwartet. Dabei soll er feststellen können, was diese spezielle Arbeit ausmacht und inwiefern sie sich von anderen Arbeiten auf dem gleichen Gebiet unterscheidet.

Zusätzlich können ein bis zwei Bilder hinzugfügt werden, sofern sie zum wesentlichen Verständnis der Arbeit beitragen (z.B. Flussdiagramm). 

\begin{keywords}
Ca. fünf Schlagwörter zur treffenden Kennzeichnung des Dokumenteninhalts in alphabetischer Reihenfolge. Die Schlagwörter dienen als erste schnelle Entscheidungshilfe ob die Arbeit für den Leser inhaltlich relevant ist.\\
Z.B. Finite Elemente Methode, Faserverbundwerkstoffe, Stabilitätsverhalten, Zugversuch, Bruchmechanik.
\end{keywords}